% Adjust:  abstract / enumerate / verbatin / minted  environments

% Check "Abstract" font and 'normalfont', and add a separating *COLON*, if
%   *  matching font, or mismatching boldness
\makeatletter

\newcommand{\abstractColonLogic}{%
  \begingroup
    \normalfont
    \edef\bodyFontFamily{\f@family}%
    \edef\bodyFontSeries{\f@series}%
    \abstractnamefont
    \edef\labelFontFamily{\f@family}%
    \edef\labelFontSeries{\f@series}%
    \def\tempColon{0}%
    \ifx\bodyFontFamily\labelFontFamily \def\tempColon{1}\fi
    \ifx\bodyFontSeries\labelFontSeries \else \def\tempColon{1}\fi
    \if\tempColon 1:%
    \fi
  \endgroup
}
\makeatother

% Format "Abstract: " as leading word (no colon if "Abstract") in different font
\renewenvironment{abstract}
    {\begin{quotation}\noindent{\abstractnamefont\abstractname\abstractColonLogic}\space\normalfont\normalcolor}
    {\end{quotation}}

% customizes ALL levels of enumerate
\usepackage{enumitem}
\setlist[enumerate]{
    label=\arabic*),
    % nosep,
    itemsep=1ex,
    parsep=0pt,
    topsep=6pt plus 2pt minus 2pt,  % restore LaTeX default
}

% prepare/setup ADJUSTED SIZE for *MINTED* and *VERBATIM* environments
\usepackage{iftex}  % need smaller size for pdflatex
\newcommand{\customsize}{\ifTUTeX 9.33\else 9\fi pt}
% \newcommand{\customsize}{\ifTUTeX{9.33}{8.88}pt}

\newcommand{\adjustedsize}{
   \fontsize{\customsize}{12pt}\selectfont  % temp baselineskip (overwritten)
   \setlength{\baselineskip}{1.2\fontdimen6\font}%
}

% adjust *verbatim* SIZE
\makeatletter
    \def\verbatim@font{\adjustedsize\ttfamily}
\makeatother

% adjust *minted* SIZE
\usepackage{minted}
\setminted{
    escapeinside=@@,  % allow manual formatting inside of minted env's
    fontsize=\adjustedsize,
}

% Remove need for \noindent after verbatim
\usepackage{noindentafter}
\NoIndentAfterEnv{verbatim}
% Doesn't work for \inputminted{lang}{file} -- that's done manually in TEX-ENGINE-MODS.tex, since this pkg could not "eat" indentations in text placed AFTER a blank line.

% These two commands are provided by TEX-ENGINE-MODS.tex for use in the main document.
\providecommand{\explicitindent}{\relax}
\providecommand{\indentpar}{\relax}
% In cases where TEX-ENGINE-MODS.tex is disabled, the commands may be left throughout the document causing compilation failures.  We "provide" them here as doing nothing, just a \relax (no indentation needs to be done if TEX-ENGINE-MODS.tex is opted out of) so that the compilar process can succeed without having to remove instances of these from the main content.
